\documentclass[11pt]{article}
\usepackage[margin=1in]{geometry}
\usepackage{amsmath,amssymb,amsthm}
\usepackage[colorlinks=true,linkcolor=blue,citecolor=blue,urlcolor=blue]{hyperref}

\newtheorem{theorem}{Theorem}
\newtheorem{proposition}[theorem]{Proposition}
\newtheorem{lemma}[theorem]{Lemma}
\theoremstyle{remark}
\newtheorem{remark}[theorem]{Remark}

\newcommand{\floor}[1]{\left\lfloor #1\right\rfloor}
\newcommand{\lcm}{\operatorname{lcm}}

\title{A charge addendum for Erd\H{o}s Problem 488: primitive quadruples}
\author{Sandbox note (Claude + Codex collaboration)}
\date{7 July 2026}

\begin{document}
\maketitle

\begin{abstract}
This addendum extends the elementary charge argument for Erd\H{o}s Problem 488
from primitive triples to primitive quadruples.  The claim is local and
unaudited: it should be read as a proof candidate pending adversarial, human, and
literature review; see \texttt{REFEREE\_QUADRUPLES.md} and
\texttt{audit\_quadruple\_charge.py} for a first local audit.  The key observation is that, for four generators, two
positive charges suffice once the triple and quadruple intersection terms are
kept.  We then show every primitive quadruple has at least two positive charges.
\end{abstract}

\section{Setup}

Let \(P=\{a,b,c,d\}\) be primitive with \(a<b<c<d\), and let
\[
  S={1\over a}+{1\over b}+{1\over c}+{1\over d}.
\]
For \(e\in P\), define
\[
  X_e(n)=\floor{n/e}-\sum_{f\in P,\ f\ne e}\floor{n/\lcm(e,f)}.
\]
Call \(e\) good if
\[
  \sum_{f\in P,\ f\ne e}{\gcd(e,f)\over f}<1. \tag{G_e}
\]
If \(e\) is good and \(n\ge d\), then \(X_e(n)\ge1\): writing
\(t=\floor{n/e}\ge1\) and using
\(\floor{n/\lcm(e,f)}=\floor{t/(f/\gcd(e,f))}\), we have
\[
X_e(n)\ge t\left(1-\sum_{f\ne e}{\gcd(e,f)\over f}\right)>0,
\]
and \(X_e(n)\) is an integer.

Write
\[
 s(n)=\sum_{e\in P}\floor{n/e},\qquad
 P_2(n)=\sum_{\{e,f\}\subset P}\floor{n/\lcm(e,f)},
\]
\[
 T_3(n)=\sum_{\{e,f,g\}\subset P}\floor{n/\lcm(e,f,g)},\qquad
 T_4(n)=\floor{n/\lcm(a,b,c,d)}.
\]
Exact inclusion-exclusion gives
\[
  B(n)=s(n)-P_2(n)+T_3(n)-T_4(n).
\]
Also \(2s(n)>nS+s(n)-4\).  Hence it is enough to prove
\[
  H(n):=s(n)-2P_2(n)+2T_3(n)-2T_4(n)\ge4,
\]
because then
\[
  2B(n)>nS+H(n)-4\ge nS.
\]

\section{Two good charges suffice}

\begin{proposition}
If at least two elements of \(P\) are good, then \(2B(n)>nS\) for every
\(n\ge d\). Consequently \(B(m)/m\le S<2B(n)/n\) for all \(m\ge1,n\ge d\).
\end{proposition}

\begin{proof}
Let \(G\subset P\) be two good elements and let \(H_0=P\setminus G\).  The good
charges contribute
\[
  \sum_{g\in G}X_g(n)\ge2.
\]
Now group the other two generators:
\[
  Y(n)=\sum_{h\in H_0}X_h(n)+2T_3(n)-2T_4(n).
\]
We claim \(Y(n)\ge2\).  Check pointwise.  For \(k\le n\), let \(p\) be the number
of elements of \(H_0\) dividing \(k\), and \(q\) the number of elements of \(G\)
dividing \(k\).  Only points with \(p>0\) contribute to \(Y\).  Their weights are
\[
\begin{array}{c|ccc}
 & q=0&q=1&q=2\\ \hline
p=1&1&0&1\\
p=2&0&0&2
\end{array}
\]
so all contributions are nonnegative.  Since \(P\) is primitive and \(n\ge d\),
each of the two generators in \(H_0\) itself contributes with \(p=1,q=0\), hence
\(Y(n)\ge2\).  Therefore
\[
H(n)=\sum_{g\in G}X_g(n)+Y(n)\ge4,
\]
and the setup above gives \(2B(n)>nS\).  The union bound gives
\(B(m)/m\le S\).
\end{proof}

\section{Every primitive quadruple has two good charges}

\begin{lemma}
The least element \(a\) is good.
\end{lemma}

\begin{proof}
For \(y\in\{b,c,d\}\), put \(q_y=y/\gcd(a,y)\).  Primitivity gives \(q_y\ge3\).
If \(q_y=3\), then \(y=3g\) with \(g=\gcd(a,y)\).  Since \(y>a\), \(g>a/3\), so
\(a/g<3\).  Also \(a/g\ne1\), else \(a\mid y\).  Thus \(a/g=2\), and
\(y=3a/2\).  This can happen for at most one \(y\).  The other two terms are
therefore at most \(1/4\), so the charge at \(a\) is at most
\[
  {1\over3}+{1\over4}+{1\over4}<1.
\]
\end{proof}

\begin{lemma}
If \(b\) is bad, then \(c\) is good.
\end{lemma}

\begin{proof}
Put
\[
q={a\over\gcd(a,b)},\qquad u={c\over\gcd(b,c)},\qquad
v={d\over\gcd(b,d)}.
\]
Primitivity gives \(q\ge2\) and \(u,v\ge3\).  Badness of \(b\) says
\[
  {1\over q}+{1\over u}+{1\over v}\ge1. \tag{1}
\]
If \(q\ge3\), then (1) forces \(u=v=3\); otherwise the sum is at most
\(1/3+1/3+1/4<1\).  But \(u=3\) forces \(c=3b/2\), and \(v=3\) forces
\(d=3b/2\), contradicting \(c<d\).  Hence \(q=2\), so \(b/a=w/2\) with \(w\)
odd, and
\[
  {1\over u}+{1\over v}\ge {1\over2}. \tag{2}
\]

Write the reduced fractions \(c/b=u/r\) and \(d/b=v/s\).  The constraints
\(2\le r<u\), \(2\le s<v\), \(c<d\), and (2) leave only five possibilities:
\[
(c/b,d/b)\in
\left\{\left({3\over2},{5\over3}\right),
\left({3\over2},{5\over2}\right),
\left({4\over3},{3\over2}\right),
\left({5\over4},{3\over2}\right),
\left({6\over5},{3\over2}\right)\right\}.
\]
Indeed, if neither \(u\) nor \(v\) is \(3\), then (2) forces \(u=v=4\), and
both ratios equal \(4/3\), impossible.  If \(u=3\), then \(c/b=3/2\) and
\(v\le6\); the order \(d/b>3/2\) leaves only \(5/3\) and \(5/2\).  If \(v=3\),
then \(d/b=3/2\), and the order \(c/b<3/2\) leaves only \(4/3,5/4,6/5\).

In these five cases the charge at \(c\) is bounded respectively by
\[
{1\over4}+{1\over2}+{1\over10},\quad
{1\over4}+{1\over2}+{1\over5},\quad
{1\over3}+{1\over3}+{1\over9},\quad
{1\over8}+{1\over4}+{1\over6},\quad
{1\over5}+{1\over5}+{1\over5},
\]
all of which are \(<1\).  Thus \(c\) is good.
\end{proof}

\begin{theorem}
For every primitive quadruple \(P=\{a,b,c,d\}\), \(2B(n)>nS\) for all \(n\ge d\).
Consequently Erd\H{o}s Problem 488 holds for every finite set whose primitive
core has size at most \(4\).
\end{theorem}

\begin{proof}
The least element \(a\) is good.  If \(b\) is good, then \(a,b\) are two good
charges.  If \(b\) is bad, then \(c\) is good, so \(a,c\) are two good charges.
The proposition applies.  Passing from a set to its primitive core preserves
\(B\) and does not increase the maximum element.
\end{proof}

\begin{remark}
This addendum is not yet priority- or literature-audited.  It should not be
presented publicly as a new result until the proof has survived adversarial and
human review.
\end{remark}

\end{document}
